\begin{center}[7.]\par\medskip\textit{The Law of Induction.}\par\medskip
{\small (Found 1835. January 23. At $7^{\mathrm u}$ in the morning, before getting up.)}
\end{center}
I.\quad The current-generating force produced at a point $P$ by a conductor element $\gamma$ whose distance from $P$ is $=r$ is, during the time element $\mathrm dt$, the difference of the two values of $\frac{\gamma}{r}$ corresponding to the instants $t$ and $t+\mathrm dt$, divided by $\mathrm dt$. Here $\gamma$ must be taken into account in both magnitude and direction. The expression
\[
-\frac{\mathrm d\dfrac{\gamma}{r}}{\mathrm dt}
\]
states this briefly and clearly.

\SourceLeaf{616}II.\quad The strength of a generated current is
\[
=\frac{\int p\,\mathrm ds}{\theta}
\]
where $p$ is the current-generating force in each element $\mathrm ds$ of the conductor and $\theta$ is the total resistance.

\begin{center}[8.]\end{center}
Let $s,s'$ be two closed conductors, $r$ the distance between two points in $s$ and $s'$, $r=\rho\rho$, and $\theta$ the total resistance in $s'$.

The two fundamental laws are as follows.

I.\quad If galvanic currents in $s$ and $s'$ have intensities $\varepsilon,\varepsilon'$, taken as positive when they run through the conductors in the direction chosen for the elements $\mathrm ds,\mathrm ds'$, then the repulsive force of the elements is
\[
=+\frac{\varepsilon\varepsilon'}{\rho}\,.\,\frac{\mathrm d\mathrm d\rho}{\mathrm ds\,.\,\mathrm ds'}\,.\,\mathrm ds\,.\,\mathrm ds'
\]
or, if $\mathrm d,\mathrm d'$ denote the partial differentiations with respect to the two currents,
\[
=+\frac{\varepsilon\varepsilon'}{\rho}\,.\,\mathrm d\mathrm d'\rho
\]
This force acts along the straight line $r$.

\vspace{3mm}
II.\quad If, during the very short time $\delta t$, the current arises in $s$, this is accompanied at every point by a current-generating force of the kind noted above. Its measure for the element $\mathrm ds'$ is
\[
=-\frac{\varepsilon\,\mathrm ds\,.\,\mathrm ds'}{\delta t\,.\,r}\cos u
\]
where $u$ denotes the angle between the directions $\mathrm ds,\mathrm ds'$.

If $z$ denotes the projection of $r$ onto the direction of $\mathrm ds'$, then $\mathrm ds\,.\,\cos u=\mathrm dz$, so the formula becomes
\[
=-\frac{\varepsilon\,\mathrm dz}{\delta t\,.\,r}\,.\,\mathrm ds'
\]
or the total current-generating force induced by $is$ in $\mathrm ds'$ is
\SourceLeaf{617}\[
=-\frac{\varepsilon\,.\,\mathrm ds'}{\delta t}\int\frac{\mathrm dz}{r}
\]

Since $\frac{\mathrm dz}{r}-\frac{z\,\mathrm dr}{rr}$ is an exact differential and its integral along the closed conductor $s$ is therefore $=0$, the integral above is also
\[
=-\frac{\varepsilon\,\mathrm ds'}{\delta t}\int\frac{z\,\mathrm dr}{rr}
\]
or, since $-\frac{z}{r}\,.\,\mathrm ds'=\mathrm d'r$,
\[
=+\frac{\varepsilon}{\delta t}\int\frac{\mathrm dr\,.\,\mathrm d'r}{r}=+\frac{4\varepsilon}{\delta t}\int\mathrm d\rho\,.\,\mathrm d'\rho
\]
so the total current-generating force is $=\frac{4\varepsilon}{\delta t}\int\int'\mathrm d\rho\,.\,\mathrm d'\rho$. Consequently, the current in $s'$ during the time $\delta t$ has intensity
\[
\frac{4\varepsilon}{\theta\delta t}\,.\,\int\int'\mathrm d\rho\,.\,\mathrm d'\rho
\]
for which one can evidently also write
\[
-\frac{4\varepsilon}{\theta\delta t}\,.\,\int\int'\rho\,\mathrm d\mathrm d'\rho
\]
Indeed, $\int\,.\,\int'\mathrm d\rho\,.\,\mathrm d'\rho=\int\rho\,\mathrm d\rho-\int\int'\rho\,\mathrm d\mathrm d'\rho$ and $\int\rho\,\mathrm d\rho=0$ since the integral is taken along the entire current line $s$.

